\subNsection{1}{Turing Maschine}{Author}

Die Turing-Maschine ist als eigenständiger Themenbereich innerhalb des Moduls
\glqq Grundlagen der Informatik\grqq{} umgesetzt. Der Bereich besteht nicht nur
aus einem Simulator, sondern aus einem vollständigen Lern-Workflow mit
Übersicht, Übungen, Beispielen und einer freien Simulation. Die Git-Historie
zeigt dabei eine schrittweise Entwicklung: Zuerst wurden Datenmodell,
Simulator-Schnittstelle und Debugger aufgebaut, danach folgten die editierbare
Konfiguration, die responsive Banddarstellung, Beispiele und Übungen,
Internationalisierung, Import/Export, automatische Speicherung und zuletzt
weitere Tests sowie die Integration in gemeinsam genutzte Komponenten.

\subNsection{2}{Entscheidungen in der Architektur}{Author}

\begin{itemize}
  \item Die Implementierung trennt die Turing-Maschine konsequent in
  \texttt{domain}, \texttt{application}, \texttt{presentation} und
  \texttt{adapters}. Dadurch liegen fachliche Regeln, React-Zustand,
  UI-Komponenten und die Abbildung auf generische Tabellenkomponenten in
  getrennten Schichten.

  \item Das statische Maschinenprogramm und der Laufzeitzustand sind getrennte
  Modelle. \texttt{TMProgram} beschreibt Zustände, Alphabete, Blank-Symbol,
  Startzustand, Haltezustände, Übergänge und Anfangsband. \texttt{TMRuntime}
  beschreibt aktuellen Zustand, Kopfposition, Bandinhalt, Schrittzähler und
  Halteinformation.

  \item Die eigentliche Simulation liegt in reinen Funktionen wie
  \texttt{createInitialRuntime}, \texttt{readCurrentSymbol} und
  \texttt{stepRuntime}. Diese Funktionen hängen nicht von React, i18n oder der
  Darstellung ab und sind deshalb gut isoliert testbar.

  \item Fehlende Übergänge werden fachlich als normales Halten modelliert und
  nicht als technischer Fehler. Ein fehlender Übergang erzeugt einen
  \texttt{haltReason} vom Typ \texttt{missing\_transition}; ein regulärer
  Haltezustand erzeugt einen \texttt{haltReason} vom Typ
  \texttt{halt\_state}. Dadurch kann die UI sauber zwischen einem erwartbaren
  Verhalten einer partiell definierten Turing-Maschine und echten
  Implementierungsfehlern unterscheiden.

  \item Die Validierung prüft Struktur und Konsistenz, aber nicht die
  Vollständigkeit der Übergangsfunktion. Es wird unter anderem geprüft, ob
  Start- und Haltezustände existieren, Symbole nur aus einzelnen Zeichen
  bestehen, das Blank-Symbol gültig ist und Übergänge auf existierende
  Zustände und Bandsymbole verweisen. Unvollständige Tabellen bleiben dagegen
  erlaubt, weil sie didaktisch und fachlich sinnvoll sind.

  \item Für die Simulation wird eine Klasse \texttt{TuringMachineSimulator}
  verwendet, die den gemeinsamen \texttt{ISimulator}-Vertrag der Anwendung
  erfüllt. Intern nutzt sie aber weiterhin die reinen Domain-Funktionen. Damit
  bleibt die Turing-Maschine kompatibel zur allgemeinen Simulationsengine,
  ohne die Fachlogik in Hooks oder UI-Komponenten zu verschieben.

  \item Der Simulator speichert eine Historie der Laufzeitzustände. Dadurch
  kann die UI neben \glqq Schritt vorwärts\grqq{} auch \glqq Schritt zurück\grqq{}
  anbieten. Die Historie liegt in der Simulator-Anwendungsschicht und bleibt
  von der Darstellung unabhängig.

  \item Die Editorlogik ist in \texttt{tmEditorReducer} und
  \texttt{useTMConfigEditor} ausgelagert. Der Reducer kümmert sich um
  Operationen wie Zustände hinzufügen, Zustände umbenennen, Bandsymbole
  umbenennen, Übergänge aktualisieren und Haltezustände verwalten. Wichtig ist
  dabei, dass Referenzen mitgeführt werden: Wird ein Zustand umbenannt, werden
  auch Startzustand, Haltezustände und Zielzustände in Übergängen angepasst.

  \item Die Übergangstabelle ist intern als verschachtelte Struktur aus
  Lesesymbol und Zustand modelliert. Für die Darstellung wird sie über
  \texttt{tmTableAdapter} in Header, Zeilen und aktive Tabellenzelle
  übersetzt. Diese Adapterentscheidung verhindert, dass die Domain-Schicht
  Tabellen- oder UI-Wissen enthält.

  \item Eingabealphabet und Bandalphabet werden nicht vollständig manuell
  gepflegt. Das Eingabealphabet wird aus dem Anfangsband abgeleitet; das
  Bandalphabet wird aus vorhandenem Bandalphabet, Anfangsband, Blank-Symbol,
  Lesesymbolen und Schreibsymbolen abgeleitet. Dadurch bleibt die
  Konfiguration für Lernende einfacher, während die Domain dennoch beide
  Konzepte kennt.

  \item Der Import und Export ist als eigener Transfer-Codec umgesetzt. Die
  aktuelle Architektur verwendet eine Bitstream-basierte, Base64-URL-sichere
  Kodierung mit Symbolregistern, Roh-UTF-8-Fallback und verschiedenen
  Übergangstabellen-Kodierungen. Die Historie zeigt, dass dies aus einem
  einfachen Base64-Proof-of-Concept zu einer kompakten, robusteren
  Speicherform weiterentwickelt wurde.

  \item Für gespeicherte Aufgabenstände wird der gemeinsame
  Training-Center-Speicher genutzt. Beispiele werden bewusst nicht dauerhaft
  überschrieben, während freie Simulationen und Übungen eigene
  Speicherkennungen verwenden. Dadurch bleiben eingebaute Beispiele
  reproduzierbar und Übungsfortschritte lokal erhalten.

  \item Die Übungsprüfung ist von der sichtbaren Simulation getrennt. Für jeden
  Validierungsfall wird ein Simulator mit dem jeweiligen Eingabewort gestartet,
  bis zum Halten oder bis zu einer Schrittgrenze ausgeführt und dann mit der
  erwarteten Ausgabe verglichen. Zusätzlich reicht eine passende Ausgabe nicht:
  Die Maschine muss in einem gültigen Haltezustand enden.
\end{itemize}

\subNsection{2}{Entscheidungen im UI}{Author}

\begin{itemize}
  \item Die Turing-Maschine nutzt den gemeinsamen Themenaufbau der Plattform:
  Kopfbereich mit Titel und Beschreibung, Tabs für Übersicht, Übungen,
  Beispiele und Simulation sowie die globale Sidebar mit Sprach- und
  Theme-Umschaltung.

  \item Die Simulation startet im Konfigurationsmodus. Lernende sehen zuerst
  Anfangsbandwort, Startzustand, Übergangstabelle, Haltezustände,
  Eingabealphabet, Bandalphabet sowie Import/Export. Erst nach dem Anwenden
  der Konfiguration wechselt die Ansicht in den Laufzeitmodus.

  \item Die Übergangstabelle ist der zentrale Editor. Zeilen stehen für
  Bandsymbole, Spalten für Nicht-Haltezustände. Haltezustände werden bewusst
  nicht als editierbare Übergangsspalten angezeigt, weil von ihnen aus keine
  weiteren Übergänge erwartet werden.

  \item Übergänge werden in der Form
  \texttt{(Next State, Write Symbol, Movement)} eingegeben. Das UI erlaubt
  leere Zellen und behandelt sie als nicht definierte Übergänge. Damit passt
  die Darstellung zur Architekturentscheidung, partielle Übergangsfunktionen
  zuzulassen.

  \item Fehler in einzelnen Tabellenzellen blockieren den Editor gezielt. Bei
  ungültigen Symbolen oder ungültigem Übergangsformat bleibt die betroffene
  Zelle editierbar, andere Änderungen und der Start der Simulation werden
  deaktiviert. So wird verhindert, dass ein halbvalider Tabellenzustand
  unbemerkt übernommen wird.

  \item Das Blank-Symbol wird sichtbar als \(\diamond\) dargestellt. Für die
  Eingabe ist zusätzlich die Tastaturform \texttt{<>} vorgesehen, die in das
  Blank-Symbol umgewandelt wird. Das reduziert Eingabehürden, ohne in der
  Darstellung ein uneindeutiges Leerzeichen verwenden zu müssen.

  \item Die Laufzeitansicht besteht aus Band, Steuerleiste, Statuskarte und
  Übergangstabelle. Das Band zeigt ein Fenster um die Kopfposition mit
  Indexbeschriftungen und einer Pfeilmarkierung für den aktuellen Kopf. Die
  Tabellenzelle des aktuellen Zustands und Symbols kann über den Adapter als
  aktive Zelle markiert werden.

  \item Das Band ist responsiv ausgelegt. Die Anzahl sichtbarer Bandzellen wird
  aus der verfügbaren Breite berechnet; Animationen werden nur für kleine
  Verschiebungen und sinnvolle Geschwindigkeiten ausgeführt. Dadurch bleibt
  die Darstellung auch bei kleinen Viewports und schnellen Simulationen
  beherrschbar.

  \item Die Steuerung verwendet die gemeinsamen Simulator-Controls der
  Plattform: Play, Pause, Schritt, Schritt zurück, Reset und
  Geschwindigkeitsregler. Damit verhält sich die Turing-Maschine ähnlich wie
  andere Simulatoren der Anwendung.

  \item Die Statuskarte zeigt aktuellen Zustand, aktuelles Symbol,
  Schrittzahl, Haltestatus und bei fehlenden Übergängen eine fachliche
  Erklärung. Bei der Live-Prüfung hielt die Maschine ohne definierten Übergang
  mit der Meldung, dass für den aktuellen Zustand und das aktuelle Symbol kein
  Übergang definiert ist.

  \item Übungen erweitern die Simulation um die Aktionen
  \glqq Lösung laden\grqq{} und \glqq Lösung prüfen\grqq{}. Das Prüfergebnis
  wird in einem Dialog als Tabelle mit Eingabe, Ausgabe, erwarteter Ausgabe,
  Bestehen und gültigem Haltezustand angezeigt.

  \item Die Oberfläche ist weitgehend lokalisiert. Eine sichtbare Auffälligkeit
  in der deutschen UI ist jedoch, dass in der Übergangstabelle der
  Zeilen-Hinzufügen-Text noch als \glqq Add new row\grqq{} erscheint. Das
  deutet auf einen fehlenden oder nicht angebundenen Übersetzungseintrag in
  einer gemeinsam genutzten Tabellenkomponente hin.
\end{itemize}

\subNsection{2}{Andere Entscheidungen}{Author}

\begin{itemize}
  \item Die Beispiele und Übungen sind als Szenarien modelliert. Beispiele
  enthalten ein fertiges Programm; Übungen enthalten zusätzlich ein
  Starterprogramm, ein Lösungsprogramm, eine Speicher-ID und mehrere
  Validierungsfälle. Damit kann derselbe Simulator für freie Arbeit, Beispiele
  und bewertbare Übungen genutzt werden.

  \item Die enthaltenen Beispiele und Übungen decken verschiedene typische
  Aufgaben ab, etwa Bit-Invertierung, unäres und binäres Inkrementieren,
  Löschen des Bandes, Suchen von Mustern, Duplizieren von Nullen,
  Binäraddition und einfache Worterkennung. Die Übungen sind nach
  Schwierigkeitsgrad sortiert.

  \item Die Testsuite ist für diesen Bereich vergleichsweise breit. Es gibt
  Unit-Tests für Domain-Logik, Formatierung, Programmvalidierung,
  Transfer-Codec, Editor-Reducer, Tabellenadapter, Bandlogik,
  Lösungskontrolle und Simulator-Hook. Die E2E-Tests prüfen zusätzlich die
  sichtbare Bedienung, Beispiel- und Übungsnavigation, Local-Storage-Verhalten,
  Import/Export, Fehlervalidierung und das Erstellen einer Bit-Inverter-Maschine
  von Grund auf.

  \item Die Git-Historie legt nahe, dass die Architektur bewusst nachträglich
  stabilisiert wurde. Nach einem ersten Datenmodell und Simulator folgten
  größere Refactorings nach einer Architekturplanung, danach responsive
  Verbesserungen, Internationalisierung, Test-IDs, Übungsintegration und
  persistente Speicherung. Das spricht für eine iterative Entwicklung, bei der
  die UI zunächst funktionsfähig gemacht und anschließend stärker in Schichten
  getrennt wurde.

  \item Besonders sinnvoll ist die Entscheidung, \glqq fehlender Übergang\grqq{}
  nicht als Programmfehler zu behandeln. Für Lernende ist eine unvollständige
  Turing-Maschine ein normaler Zwischenzustand beim Konstruieren. Gleichzeitig
  kann die Übungsprüfung streng bleiben, weil sie nur dann besteht, wenn die
  Maschine nicht nur die richtige Ausgabe produziert, sondern auch in einem
  vorgesehenen Haltezustand endet.

  \item Ein technischer Zielkonflikt liegt im Transferformat. Die kompakte
  Bitstream-Kodierung ist gut für kurze Exportstrings und lokale Speicherung,
  erschwert aber die manuelle Lesbarkeit. Diese Entscheidung ist plausibel,
  weil der Import/Export im UI als Kopier- und Austauschfunktion gedacht ist
  und nicht als primäres Bearbeitungsformat.

  \item Ein weiterer Zielkonflikt liegt in der abgeleiteten Alphabetpflege.
  Sie reduziert Eingabeaufwand und verhindert viele Inkonsistenzen, kann aber
  für fortgeschrittene theoretische Diskussionen weniger explizit sein als ein
  separat editierbares Eingabealphabet. Die aktuelle Lösung wirkt für die
  Lernplattform angemessen, weil sie den Editor einfacher hält.

  \item Aus Wartungssicht ist positiv, dass die meisten UI-Elemente stabile
  \texttt{data-testid}-Attribute besitzen. Das erleichtert E2E-Tests und war
  auch bei der Live-Inspektion hilfreich. Gleichzeitig sollten kleine
  Lokalisierungslücken in gemeinsam genutzten Komponenten weiter beobachtet
  werden, damit deutsche und englische Oberfläche konsistent bleiben.
\end{itemize}
